\section{Evaluation and Conclusions}

The analyses of results obtained using the methodology proposed in this study show great
promise for the future of conservational biomonitoring.
It is shown that application of current state-of-the-art monocular depth estimation
frameworks in conjunction with modern distance sampling can produce estimates of population
density and abundance on the order of those derived in the convention manner.
Most importantly though, the time and labour commitments necessitated by this methodology
are orders of magnitude lower.
Manual distance estimation across a dataset of the size used in this study may take weeks,
even months, depending on available resources.
The pipeline explored in this study achieves this in hours.
Nevertheless, methodology must continue to be adapted with the emergence of novel models
and technologies before it can become convention.

\subsection{Roadblocks and Bottlenecks}

This study has found that, when outsourcing distance estimation to an automated pipeline,
resulting estimates for density and abundance consistently fall below those using manual
distances.
However, the greatest barrier that must be overcome here is not error in distance estimates,
although this remains important, but the endeavour of detection itself.
This is certainly true for, at least, this dataset.
It is inevitable that automated detection will result in misses, but the frequency of these
is significant and noticeably impacts estimated density since misses occur over a range
of distances, diminishing the effectiveness of the fitted detection function to account
for them.

Fundamentally, detections are missed because of a lack of contextual cues in detection
frames.
Using the manual approach, the human annotator has additional context from subsequent and
previous video frames.
Future studies may attempt to replicate this idea in the distance estimation pipeline.
As it stands, the pipeline of this study uses Mega Detector to identify detections
in extracted frames.
Suppose instead that animal bounding boxes are tracked, continuously, throughout the
duration of the camera track video using, for example, YOLOv8.\cite{YOLOv8}
The positions of these bounding boxes may then be sampled at intervals corresponding to the
extracted frames and then processed with the depth data in an equivalent manner.
This method is likely to lower the frequency of misses but by how much is yet to be seen.

Somewhat of a bottleneck in this methodology is the preparation of the DPT calibration
material (i.e., calibration frames and binary masks).
The time requirement for preparing these is certainly a huge improvement over time saved
manually annotating detection distances, but is still significant, most especially for
binary mask preparation which took approximately a week to complete.
As discussed, the automated masking method in Section~\ref{subsubsec:mask_creation} yields
inconsistent results, highlighting the need for refinement here.
The custom-built calibration frame extractor software, as seen in Section~\ref{fig:frame_extractor},
proved to be of great success however, allowing for all raw calibration frames to be extracted at a
significantly increased rate.


available compute

Time saved using frame extractor, preserves exact pixels

Practitioners are not concerned with manual truncation, binning and model fitting for density
estimation as these require minimal effort.

Limitations of the environment\\
Missed detections = low density, aka mega detector is shite\\
Future work, use detector model that tracks animal bounding boxes through video to extract
bounding boxes, gives more context, less misses, better density estimates.

\subsection{Optimising Estimates}

When applied to this dataset, DPT outperforms Depth Anything with regard absolute error of
distance estimates.
Additionally, while both models over-predict distance to some degree, the skew is much more
neutral in the case of DPT, in comparison to Depth Anything which significantly over-predicts.
As explained in Section~\ref{subsubsec:configuration_effects}, this leads to DPT being the
better choice of distance estimation model for achieving density and abundance results that are
most reflective of the ground-truth.
At a high level, this makes sense given that DPT output is specifically calibrated
to each individual camera location and is therefore a better description of the environment.
Future studies may show interesting results if a similar reference frame style calibration
is to be applied to the Depth Anything.
This has the potential to produce better estimates of density and abundance than DPT,
given Depth Anything's reasonable performance without calibration and its ability better
capture fine detail in depth contrast than DPT\@.

In Section~\ref{subsubsec:var_cal}, it was demonstrated that, when DPT is calibrated using
only the closest and furthest reference frame, the estimated distances of close pixel are
underestimated, while those of far pixels are overestimated.
In an attempt to explain this, a hypothesis was proposed theorising that the true relationship
between DPT outputted disparity and metric disparity follows a non-linear mapping.
To the best of the knowledge of the author, the literature contains no studies that map
DPT disparity to metric disparity.
As such, this relationship should be investigated.
Besides validation this hypothesis, evidence of non-linearity would give insight into improved
methods of DPT calibration for more robust predictions.

More calibration frames


All in all, considering performance as well as necessitated setup requirements, it should be understood that
both distance estimation models have their own unique benefits and draw-backs.






Depth estimation point consistency (the animal itself has depth)\\
Error in reference videos (person has depth, not always standing up straight)


Better calibration (polynomial/not just linear, more calibration frames – across frame width,
shorter distance intervals)\\
Future work should determine true disparity relationship, indicated by the reduced calibration

Calibration and detection frame image resolution?


\subsection{Closing Remarks}

Summary of work

Novel

Improvement over fully manual approach?

Ballpark density and abundance estimates\\
MDE and detection models are always improving and estimates will only improve

Significant time improvement

Distance pipeline requires work to lower rate of missed detections which skews density estimates

